\uuid{pK7n}
\titre{Points critiques d'une fonction de deux variables}
\niveau{L2}
\module{Analyse}
\chapitre{Fonction de plusieurs variables}
\sousChapitre{Extremums locaux}
\theme{Hessienne, point selle, minimum local}
\auteur{}
\datecreate{2026-04-08}
\organisation{}
\difficulte{2}

\contenu{
	\texte{
		Soit $f : \mathbb{R}^2 \to \mathbb{R}$ définie par
		\[
		f(x, y) = x^3 - 3x + y^2.
		\]
	}
	\begin{enumerate}
		
		\item
		\question{Déterminer les points critiques de $f$.}
		\indication{Résoudre le système $\nabla f(x,y) = (0,0)$.}
		\reponse{
			On calcule les dérivées partielles :
			\[
			\frac{\partial f}{\partial x}(x,y) = 3x^2 - 3, \qquad
			\frac{\partial f}{\partial y}(x,y) = 2y.
			\]
			Le système $\nabla f = 0$ donne $x = \pm 1$ et $y = 0$.
			Les points critiques sont donc $A = (1, 0)$ et $B = (-1, 0)$.
		}
		
		\item
		\question{Pour chaque point critique, déterminer sa nature à l'aide de la matrice hessienne.}
		\indication{Calculer $\det H_f$ et le signe de $\frac{\partial^2 f}{\partial x^2}$ en chaque point critique.}
		\reponse{
			La matrice hessienne de $f$ est
			\[
			H_f(x,y) = \begin{pmatrix} 6x & 0 \\ 0 & 2 \end{pmatrix}.
			\]
			\begin{itemize}
				\item En $A = (1, 0)$ : $\det H_f = 12 > 0$ et $\frac{\partial^2 f}{\partial x^2}(1,0) = 6 > 0$,
				donc $A$ est un \textbf{minimum local}.
				\item En $B = (-1, 0)$ : $\det H_f = -12 < 0$,
				donc $B$ est un \textbf{point selle} (ni maximum ni minimum local).
			\end{itemize}
		}
		
	\end{enumerate}
}